\begin{table*}[htbp]
\centering
\begin{threeparttable}
% AAAI allows reducing font size to 9pt. \small is perfect for this.
\small 
\setlength{\tabcolsep}{4pt}
\begin{tabular}{@{}ll p{4cm} p{7cm}@{}}
\toprule
\textbf{Dataset} & \textbf{Attribute} & 
\multicolumn{1}{c}{\parbox{4.5cm}{\centering\textbf{Literature Ground Truth}\\[2pt](Priv. / Unpriv. Groups)}} & 
\multicolumn{1}{c}{\parbox{6.5cm}{\centering\textbf{LLM Consensus}\\[2pt](Tier: Priv. / Unpriv. Groups)}} \\
\midrule

\multirow{6}{*}{Law School\tnote{a}}
 & \textbf{race} & White / Non-White & T1:  White / Non-White \\
  & \textbf{gender} & Male / Female & T1: Male / Female \\
 & \textbf{family income} & - & T2: Higher Income / Lower Income \\
 & \textbf{LSAT score} & - & T2: Higher LSAT Scores / Lower LSAT Scores \\
 & \textbf{UGPA score} & - & T2: Higher UGPA Scores / Lower UGPA Scores \\
 & \textbf{law school tier} & - & T2: Higher Tier (Elite) / Lower Tier \\
\midrule
\multirow{3}{*}{COMPAS\tnote{b} }
 & \textbf{race} & White / Black & T1: White / Non-white \\
 & \textbf{gender} & Male / Female & T1: Male / Female \\
& \textbf{age} & - & T2: Older ($>$ 45) / Younger ($<$ 25)\\
\midrule
\multirow{5}{*}{Dutch Census\tnote{c}}
 & \textbf{gender} & Male / Female & T1: Male / Female \\
 & \textbf{county birth} & - & T1: Dutch-born / Foreign-born\\
 & \textbf{citizenship} & - & T1: Dutch citizens / Non-Dutch citizens \\
 & \textbf{age} & - & T2: Middle (25--54) / Young ($<$25) \& Older ($>$55) \\
 & \textbf{marital status} & - & T2: Married / Non-married \\
\midrule
\multirow{4}{*}{Bank Marketing\tnote{d}}
 & \textbf{marital status} & Married / Non-married & T1: Married / Non-married \\
 & \textbf{age} & Middle (25--60) / Young ($<$25) \& Older ($>$60) & T1: Middle (30--55) / Young ($<$30) \& Older ($>$60) \\
& \textbf{job type} & - & T2: Management / Blue-collar \\
& \textbf{education} & - & T2: Higher education / Lower education\\
\bottomrule
\end{tabular}

\caption{Validation of LLM-derived fairness context against academic literature~\cite{le2022survey, coraglia2024evaluating, kartalmitigating}.  Due to high cross-model consistency, we present a single \textbf{LLM Consensus} column. T1 denotes strongly protected attributes, and T2 denotes context-dependent. Minor variations are noted below.}
\label{tab:llm_validation}

\begin{tablenotes}

    \item[a] \textbf{Law School:} The LLMs reached a consensus on Tier 2, noting that while UGPA and LSAT scores appear merit-based, they are closely tied to socioeconomic privilege. The models reasoned that unequal access to resources such as expensive test preparation makes these metrics potent proxies for class and race, risking the reinforcement of systemic inequalities.

    \item[b] \textbf{COMPAS:} All models agreed on Tier 2 for age, but diverged in defining unprivileged groups. Gemini identified younger individuals as disadvantaged, warning that reliance on the ``age-crime curve'' risks creating a self-fulfilling prophecy based on an immutable trait. Claude, meanwhile, considered both younger and older individuals as unprivileged, surfacing a nuanced ethical debate on the specific nature of age-based vulnerability in criminal risk assessment.

    \item[c] \textbf{Dutch Census:} Minor disagreements revealed distinct model tendencies. GPT-4.1 classified \texttt{citizenship} as a structural advantage rather than a directly ascriptive protected attribute in the Meritocratic framework, while all three models agreed on Tier 1 for \texttt{country\_birth} across all frameworks. GPT-4.1 also excluded marital status from protected attributes in the Anti-Classification framework, assigning it a mixed/context-dependent role in the Meritocratic framework. Gemini and Claude provided more granular reasoning, suggesting that privileged/unprivileged groups for marital status should be stratified by gender.

    \item[d] \textbf{Bank Marketing:} Gemini and Claude classified age and marital status as Tier 1, citing EU and Portuguese anti-discrimination laws, and explicitly referenced the 2008 financial crisis, noting that age-based risk assessment was particularly problematic in this historical context. GPT-4.1 classified age as Tier 1 but did not treat marital status as a directly protected attribute under Portuguese anti-classification law (though it recognized it as an ascriptive attribute in the Meritocratic framework). Job type and education were assigned Tier 2 by all models, as they identified risks to financial inclusion for blue-collar populations, given that the dataset was collected during a period of economic instability.

\end{tablenotes}
\end{threeparttable}
\end{table*}
